\pg{441} produziu e formou. Logo o atheismo \emph{theoretico-negativo} é impossivel nos homens \emph{adultos} \footnote{A historia contém exemplos de homens que, isolados, desde crianças, de toda a convivencia social, chegaram, só pela força da razão, ao conhecimento de uma causa productora do mundo.}.

20. \textbf{O atheismo theoretico-positivo é impossivel.} --- E' impossivel que o homem esteja \emph{séria} e \emph{sinceramente} convencido da falsidade de uma coisa, attestada pela voz da natureza. Ora a existencia de \textsc{Deus} é attestada pela voz da natureza; pois é admittida por todos os homens, em todos os tempos, em todos os logares, e só o que deriva da natureza se encontra em todos os homens, em todos os tempos, em todos os logares. Logo é impossivel que o homem esteja \emph{séria} e \emph{sinceramente} convencido de que não existe \textsc{Deus}. --- A experiencia confirma a these; pois os que se dizem \emph{atheus} são \emph{scepticos} ou \emph{pantheistas} \footnote{O nosso seculo foi chamado \emph{era do atheismo}. E' certo que alguns infelizes empregam esforços para destruir, se fosse possivel, a propria idea de \textsc{Deus}. \emph{Fouerbach} diz: « Só o homem é o verdadeiro salvador! Só o homem é o nosso \textsc{Deus}, o nosso Juiz, o nosso Redemptor ». \emph{Rénan} escreve: « \textsc{Deus}, a Providencia, a immortalidade, não passam de palavras velhas e pesadas ». A loja maçonica de Liege, em 1865, dizia: « Não só devemos collocar-nos acima de todas as religiões, mas acima de toda a crença num Deus ».\par Mas as pretendidas \emph{convicções} d'esses desvairados não passam de \emph{duvidas}. O espirito nunca poderá sinceramente convencer-se de que \textsc{Deus} não existe e, pelo menos com as suas duvidas, protestará sempre contra as estultas affirmações do coração. Já nos seus tempos \emph{Seneca} dizia: « Mentem os que dizem estar convencidos de que \textsc{Deus} não existe; pois, se ás vezes o dizem, todavia de noite, quando estão sós, começam a duvidar ».\par Estas \emph{duvidas} ácerca da existencia de \textsc{Deus} derivam da \emph{superficialidade} do pensamento e, sobretudo, da \emph{corrupção} da vontade.\par As \emph{duvidas} derivam da \emph{superficialidade} do pensamento. A este respeito, \emph{Caro} escreve: « Ha entre nós uma classe de espiritos scientificamente desorganisados, habeis para comprehender os phenomenos e as leis, nas diversas ordens da realidade, na natureza ou na historia, mas incapazes, por causa de uma inercia systematica, que se toma por uma força, de se elevar á concepção da \emph{causa}, que é a razão ultima das leis naturaes, assim como as leis são a razão dos factos » (\emph{L'idée de Dieu}, 1873, pag. 42).\par Mas as \emph{duvidas} derivam principalmente da \emph{corrupção} da vontade. Não é na \emph{intelligencia}, que o atheu nega a existencia de \textsc{Deus} ou duvida d'ella, mas é no \emph{coração}. O homem, cuja vida é pura e ordenada, nem sequer duvida da existencia de \textsc{Deus}. \emph{Rousseau} dizia: « Conservai a vossa alma num estado de desejar que exista \textsc{Deus}, e nunca duvidareis da sua existencia ». \emph{La Bruyère} escreve: « Desejava ouvir a um homem sabio, moderado, casto e justo, que \textsc{Deus} não existe ou que a nossa alma não é immortal; elle fallaria sem interesse. Mas esse homem não se encontra ».}.

21. \textbf{Consequencias do atheismo.} --- O \emph{atheismo} é a destruição da \emph{sciencia}, da \emph{sociedade} e da \emph{natureza humana}.
